\section{Supplier Competition and Policy Regimes}
\label{sec:main-results}

We now vary the exogenous number of suppliers. The first comparative static is elementary: adding an independent supplier lowers the expected minimum cost strictly, so
\begin{equation}
m_{N+1}<m_N
\quad\Longrightarrow\quad
G_{N+1}>G_N.
\label{eq:G-increase}
\end{equation}
The second comparative static comes from procurement theory. Under DRHR, expected winning supplier rent is strictly decreasing and converges to zero \citep{Li2005,XuLi2019}:
\begin{equation}
R_{N+1}<R_N,
\qquad
\lim_{N\to\infty}R_N=0.
\label{eq:R-decrease}
\end{equation}
The economic contribution comes from combining these known lower-stage effects with the government game.

\subsection{Competition--coordination divergence}

\begin{theorem}[Competition--coordination divergence]
\label{thm:divergence}
Under the baseline assumptions,
\begin{equation}
A_{N+1}^P<A_N^P
\qquad\text{and}\qquad
A_{N+1}^N>A_N^N.
\label{eq:opposite-thresholds}
\end{equation}
Thus stronger supplier competition moves the coordinated and decentralized differentiation thresholds in opposite directions.
\end{theorem}

The first inequality follows because competition raises expected real project surplus. The second follows because competition dissipates the locally captured supplier rent that rewards a region for accepting $D$. The same change in market structure therefore makes differentiation more attractive to a coordinator and less attractive to the decentralized region that must specialize on the supplier side.

Define the inefficient-duplication interval
\begin{equation}
\mathcal W_N\equiv(A_N^P,A_N^N).
\label{eq:wedgeN}
\end{equation}
For every $A\in\mathcal W_N$, Proposition \ref{prop:nash} implies that $(U,U)$ is the unique Nash equilibrium while Proposition \ref{prop:planner} implies that the coordinator strictly chooses a differentiated profile.

\begin{corollary}[Strict expansion of unique inefficient duplication]
\label{cor:wedge-expansion}
Under the baseline assumptions,
\begin{equation}
\mathcal W_N\subsetneq \mathcal W_{N+1}.
\label{eq:wedge-nesting}
\end{equation}
\end{corollary}

This result differs from a conventional statement that competition improves or worsens welfare. Conditional on differentiation, procurement becomes more efficient as $N$ rises. What expands is the set of environments in which decentralized governments fail to select the differentiated configuration in the first place.

\subsection{Competition-induced policy reversal}

The threshold result has a sharper equilibrium implication. Let $c_0$ be the lower endpoint of supplier costs. Under the maintained compact support and positive density, the minimum order statistic converges to $c_0$, and hence
\begin{equation}
\lim_{N\to\infty}G_N=v-c_0.
\label{eq:Glimit}
\end{equation}
At the same time, DRHR implies $R_N\to0$.

\begin{theorem}[Competition-induced policy reversal]
\label{thm:reversal}
Suppose
\begin{equation}
A(v-c_0)>\Delta.
\label{eq:reversal-condition}
\end{equation}
Then there exists a finite $\bar N$ such that for all $N\ge\bar N$,
\begin{equation}
\operatorname{argmax}V=\{(U,D),(D,U)\},
\qquad
NE_N=\{(U,U)\}.
\label{eq:reversal-result}
\end{equation}
\end{theorem}

Theorem \ref{thm:reversal} is the main result. If the asymptotic real surplus from a differentiated project is large enough to cover the sacrificed direct premium, a coordinator eventually prefers differentiation. Yet competition drives locally captured supplier rent toward zero, so the $D$ role eventually becomes privately unattractive to either regional government. The market can therefore approach productive efficiency within the procurement subgame while decentralized policy converges toward the wrong policy configuration.

A useful corollary starts from a market in which differentiated Nash equilibria already exist.

\begin{corollary}[Destruction of differentiated Nash equilibria]
\label{cor:destruction}
Suppose that for some $N_0$,
\begin{equation}
A\rho R_{N_0}>\Delta
\label{eq:initial-diff}
\end{equation}
so that $(U,D)$ and $(D,U)$ are pure Nash equilibria. Then there exists a finite $N_1>N_0$ such that for all sufficiently large $N\ge N_1$, those differentiated equilibria no longer exist and $(U,U)$ is the unique Nash equilibrium. If the coordinator differentiates at $N_0$, it continues to differentiate for every $N>N_0$.
\end{corollary}

This corollary gives the result an equilibrium-set interpretation. Supplier competition need not merely widen an inefficiency wedge around a fixed equilibrium. It can change the qualitative regime from one with efficient differentiated equilibria to one with unique duplicated priorities.

\subsection{Uniform illustration}

For intuition, let $F$ be uniform on $[0,1]$ and set $\rho=1$. Then
\begin{equation}
m_N=\frac{1}{N+1},
\qquad
R_N=\frac{1}{N+1},
\qquad
B_N=v-\frac{2}{N+1}.
\label{eq:uniform-objects}
\end{equation}
The thresholds are
\begin{equation}
A_N^P=\frac{\Delta(N+1)}{v(N+1)-1},
\qquad
A_N^N=\Delta(N+1).
\label{eq:uniform-thresholds}
\end{equation}
Thus the decentralized threshold grows linearly with supplier count, whereas the coordinated threshold falls toward $\Delta/v$. Figure \ref{fig:uniform-regimes} illustrates the opposite movements for $\Delta=1$ and $v=2$. The horizontal line $A=5$ shows the equilibrium-set transition directly: differentiated pure equilibria exist at low $N$, the decentralized threshold is reached at $N=4$, and for $N\ge5$ the unique Nash equilibrium is duplication, while coordinated differentiation remains optimal throughout.

\begin{figure}[ht]
\centering
\includegraphics[width=0.82\textwidth]{../figures/uniform_regime_thresholds.pdf}
\caption{Competition moves the coordinated and decentralized differentiation thresholds in opposite directions. Uniform costs on $[0,1]$, $\rho=1$, $\Delta=1$, and $v=2$. The fixed line $A=5$ illustrates the destruction of differentiated Nash equilibria as $N$ rises.}
\label{fig:uniform-regimes}
\end{figure}

The uniform case is only an illustration; Theorems \ref{thm:divergence} and \ref{thm:reversal} use the general DRHR class.
